% poster_icdm.tex -- ICDM poster.
%
% A0 portrait (841 x 1189 mm), three columns. Build with:
%     pdflatex poster_icdm.tex
%
% Numbers here are pulled from the camera-ready results in empirical-icdm/
% and simulations/sim2_se_calibration/results_icdm/, not from the originally
% submitted tables. Where the two differ it is because the submitted bootstrap
% warm-started its replicates; see the camera-ready for the correction.
\documentclass[final]{beamer}
\usepackage[size=a0,orientation=portrait,scale=1.5]{beamerposter}
\usetheme{default}
\usecolortheme{default}

\usepackage[T1]{fontenc}
\usepackage{lmodern}
\usepackage{amsmath, amssymb}
\usepackage{booktabs}
\usepackage{array}
\usepackage{xcolor}

% ---------------------------------------------------------------- palette --
\definecolor{ink}{HTML}{1A1A2E}
\definecolor{accent}{HTML}{0F4C81}
\definecolor{warm}{HTML}{B4451F}
\definecolor{panel}{HTML}{F4F1EA}
\definecolor{rule}{HTML}{C9C2B4}

\setbeamercolor{background canvas}{bg=white}
\setbeamercolor{block title}{fg=white,bg=accent}
\setbeamercolor{block body}{fg=ink,bg=panel}
\setbeamertemplate{navigation symbols}{}
\setbeamertemplate{caption}{\insertcaption}
\setbeamertemplate{itemize item}{\color{accent}$\blacktriangleright$}
\setbeamertemplate{block begin}{
  \vskip.5ex
  \begin{beamercolorbox}[leftskip=1cm,colsep*=.75ex]{block title}
    \usebeamerfont{block title}\large\bfseries\insertblocktitle
  \end{beamercolorbox}
  \vskip-.5ex
  \begin{beamercolorbox}[colsep*=.75ex,vmode]{block body}
  \vskip.5ex
}
\setbeamertemplate{block end}{\vskip.5ex\end{beamercolorbox}\vskip1.5ex}

\newcommand{\ame}{\widehat{\mathrm{AME}}}
\newcommand{\keynum}[1]{{\color{warm}\bfseries #1}}

\title{Model-Agnostic Interpretation of Machine Learning\\ via Average Marginal Effects}
\author{Jason Parker}
\institute{Southwest Airlines}
\date{}

\begin{document}
\begin{frame}[t]

% ------------------------------------------------------------------ header --
\begin{beamercolorbox}[wd=\paperwidth,colsep*=.5ex]{block title}
  \vskip1ex
  \centering
  {\Huge\bfseries Model-Agnostic Interpretation of Machine Learning}\\[.6ex]
  {\Huge\bfseries via Average Marginal Effects}\\[1.2ex]
  {\Large Jason Parker \quad\textbar\quad Southwest Airlines}\\[.6ex]
  {\large IEEE International Conference on Data Mining (ICDM) 2026}
  \vskip1ex
\end{beamercolorbox}
\vskip2ex

\begin{columns}[t]

% ================================================================ column 1 ==
\begin{column}{.31\paperwidth}

\begin{block}{The problem}
The accuracy--interpretability tradeoff is usually treated as fundamental:
either fit a linear model you can read, or fit a black box you cannot.

We argue it is an artifact of \emph{how} black boxes are read. A logistic
regression is interpretable not because it is linear but because it reports
a table of effects with standard errors. Any model can be made to do that.
\end{block}

\begin{block}{The estimand}
For a fitted prediction function $\hat f$ and feature $x_j$, the average
marginal effect is the average partial derivative over the covariate
distribution,
\[
  \mathrm{AME}_j \;=\; \mathbb{E}\!\left[\frac{\partial \hat f(\mathbf{x})}{\partial x_j}\right].
\]
This is standard in the econometrics of nonlinear models. It is well defined
for \emph{any} supervised learner, and reduces to the usual coefficient when
$\hat f$ is a generalized linear model.
\end{block}

\begin{block}{The estimator}
Analytical derivatives are replaced by centered finite differences with an
adaptive step size,
\[
  \ame_j = \frac{1}{n}\sum_{i=1}^{n}
    \frac{\hat f(\mathbf{x}_i + h_j e_j) - \hat f(\mathbf{x}_i - h_j e_j)}{2h_j},
\]
with the step size scaled to each feature and floored on the integers,
\[
  h_j =
  \begin{cases}
    \max\!\left(0.05\,\hat\sigma_j,\; \tfrac12\right)
      & x_j \text{ integer-valued},\\[.5ex]
    0.05\,\hat\sigma_j & \text{otherwise.}
  \end{cases}
\]

\begin{itemize}
  \item The integer floor matters for tree ensembles: predictions are
        piecewise constant, and a step that never crosses a split threshold
        returns a derivative of exactly zero.
  \item Binary features use the contrast
        $\hat f(\mathbf{x}^{(j,1)}) - \hat f(\mathbf{x}^{(j,0)})$ instead.
\end{itemize}
\end{block}

\begin{block}{Inference}
Standard errors come from a nonparametric bootstrap over
$(\mathbf{x}_i, y_i)$ pairs. Each replicate is \textbf{refit from a fresh
initialization} on the resample and evaluated there.

\begin{itemize}
  \item Refitting from scratch, rather than warm-starting from $\hat f$, is
        what makes the replicate spread reflect how far the fitted function
        actually moves with the data.
  \item Output is a coefficient table --- estimate, standard error,
        significance --- with no restriction on the form of $\hat f$.
\end{itemize}
\end{block}

\end{column}

% ================================================================ column 2 ==
\begin{column}{.31\paperwidth}

\begin{block}{Effects agree across model classes}
UCI Adult Income ($n = 48{,}842$), full specification. AME on
$P(\text{income} > \$50\text{k})$.

\vskip1ex
\centering
\begin{tabular}{lcccc}
\toprule
 & Logistic & Forest & XGBoost & Net \\
\midrule
education   & 0.031 & 0.029 & 0.031 & 0.028 \\
married     & 0.275 & 0.245 & 0.236 & 0.239 \\
female      & $-$0.027 & $-$0.022 & $-$0.012 & $-$0.027 \\
\midrule
McFadden $R^2$ & 0.400 & 0.450 & 0.492 & 0.455 \\
\bottomrule
\end{tabular}
\vskip1ex
\raggedright
Four learners that differ substantially in fit report the same effect of
education to within \keynum{0.3 percentage points}. This comparison is only
possible because the AME puts every model on a common scale.
\end{block}

\begin{block}{Effects in the units of the problem}
Ames Housing ($n = 1{,}460$). AME on sale price, in dollars.

\vskip1ex
\centering
\begin{tabular}{lrr}
\toprule
 & Linear & Forest \\
\midrule
Living area (/sq ft) & \$66\phantom{,000} & \$56\phantom{,000} \\
 & \small(13.06) & \small(5.25) \\
Overall quality (/pt) & \$16,465 & \$15,636 \\
 & \small(1,287) & \small(1,643) \\
\bottomrule
\end{tabular}
\vskip1ex
\raggedright
Standard errors in parentheses. The ensemble is not uniformly more precise:
it wins where the linear model must fit a nonlinear relationship with one
slope, and loses where the relationship is nearly linear.
\end{block}

\begin{block}{Why not SHAP?}
SHAP values, PDP slopes and LIME coefficients are \emph{descriptive}. There is
no standard error, no confidence interval, no hypothesis test.

They are also not on the scale of the outcome. Adult, logistic, full
specification:

\vskip1ex
\centering
\begin{tabular}{lrrr}
\toprule
 & AME & SHAP & PDP \\
\midrule
age            & 0.0027 & \keynum{$-$0.0270} & 0.0028 \\
education      & 0.0306 & \keynum{$-$0.1186} & 0.0269 \\
married        & 0.2752 & $-$0.0768 & 0.2749 \\
female         & $-$0.0269 & 0.0022 & $-$0.0270 \\
\bottomrule
\end{tabular}
\vskip1ex
\raggedright
The AME and the PDP slope agree to three decimals on three of four features
--- as they should, since both are averaging the same surface. SHAP disagrees
in \textbf{magnitude and in sign}: it reports age and education as pushing
income \emph{down} in a model whose coefficients on both are positive and
precisely estimated.

For above-ground living area in Ames, the SHAP attribution is roughly
\keynum{50$\times$} the AME for the same linear model --- not interpretable as
a marginal effect in any standard sense.
\end{block}

\end{column}

% ================================================================ column 3 ==
\begin{column}{.31\paperwidth}

\begin{block}{Simulation 1: the estimator recovers the truth}
$x_1,\dots,x_4 \sim N(0,1)$ independently; $\eta = 2x_1 + 3x_2$; $x_3, x_4$ are
noise with a true effect of zero. 500--1{,}000 Monte Carlo iterations.

\begin{itemize}
  \item Bias is negligible for every learner at every sample size on the
        noise features.
  \item Logistic and the network recover the signal effects at the
        $O(n^{-1/2})$ rate.
  \item Random forests develop a \textbf{negative bias that does not vanish}
        with $n$ --- regularization from depth capping and subsampling.
\end{itemize}
\end{block}

\begin{block}{Simulation 2: where the bootstrap fails}
Coverage of nominal 95\% intervals, linear classification design.

\vskip1ex
\centering
\begin{tabular}{lccc}
\toprule
 & $n{=}250$ & $n{=}1{,}000$ & $n{=}5{,}000$ \\
\midrule
Logistic & 0.934 & 0.930 & 0.924 \\
XGBoost  & 0.938 & 0.888 & 0.994 \\
Forest   & 0.992 & 0.978 & \keynum{0.802} \\
\bottomrule
\end{tabular}
\vskip1ex
\raggedright
\textbf{The failure is bias, not variance.} The forest's bootstrap standard
error is \emph{conservative} --- 1.2 to 1.5$\times$ the true sampling
deviation --- so the intervals are, if anything, too wide.

What breaks is that the bias stays put while the interval shrinks:

\vskip1ex
\centering
\begin{tabular}{lccc}
\toprule
Forest, $x_1$ & $n{=}250$ & $n{=}1{,}000$ & $n{=}5{,}000$ \\
\midrule
bias        & 0.013 & $-$0.002 & $-$0.016 \\
std.\ error & 0.061 & 0.023 & 0.007 \\
bias / s.e. & 0.21 & 0.08 & \keynum{2.11} \\
\bottomrule
\end{tabular}
\vskip1ex
\raggedright
At $n = 5{,}000$ the half-width of the interval (0.0145) is smaller than the
bias (0.0157). XGBoost, whose bias \emph{does} vanish with $n$, keeps nominal
coverage throughout --- so the failure is specific to learners whose
regularization bias does not shrink.
\end{block}

\begin{block}{Takeaways}
\begin{itemize}
  \item Any supervised model can report a coefficient table on the scale of
        the outcome, with inference.
  \item Effects are stable across learners that differ substantially in fit.
  \item The refitting bootstrap is valid where the learner's regularization
        bias is negligible, and \textbf{fails silently where it is not}.
  \item Open: bootstrap validity is conjectured, not proved. A companion
        journal paper establishes gradient consistency for neural networks,
        gradient boosting and Mondrian forests, and replaces the bootstrap
        with a debiased, cross-fitted, Neyman-orthogonal estimator that
        restores nominal coverage for Breiman forests.
\end{itemize}
\end{block}

\begin{block}{Software}
\centering
\texttt{marginfx} --- \texttt{pip install marginfx}\\[.5ex]
\small One call per model. scikit-learn, XGBoost, TensorFlow and PyTorch.
\end{block}

\end{column}
\end{columns}
\end{frame}
\end{document}
